\section{Переход к специализированным энергоэффективным вычислениям}

\subsection{Масштабирование электроники и изменение ограничений}

В работе Р.~Деннарда и соавторов 1974 года рассматривалось
согласованное уменьшение размеров транзисторов и рабочих
напряжений при изменении концентрации примеси~\cite{dennard1974}.
При сохранении напряжённости электрического поля в идеальной
модели сокращались задержка и мощность отдельной схемы,
а плотность мощности оставалась постоянной. Поэтому рост числа
элементов на единицу площади мог сочетаться с прежними условиями
отвода тепла. Эти соотношения описывают электрические свойства
элементов, а не только увеличение их количества на кристалле.

Рост числа транзисторов необходимо отличать от увеличения числа
одновременно активных вычислительных элементов. В исследовании Х.~Эсмаилзаде
и соавторов, опубликованном в развёрнутом виде в 2012 году,
моделировались ограничения дальнейшего масштабирования многоядерных
процессоров при заданном ограничении мощности~\cite{dark2012}.
Термин \textit{dark silicon} обозначает ограничение доли элементов
кристалла, способных одновременно работать в пределах установленного
предела мощности. Снижение напряжения питания не могло продолжаться
в прежнем темпе из-за ограничений порогового напряжения
и токов утечки. Поэтому уменьшение размеров само по себе
перестало обеспечивать постоянную плотность мощности.

При ограниченной мощности необходимо учитывать не только количество
арифметических элементов, но и объём работы, который они способны выполнить с
доступными энергетическими ресурсами. Специализация позволяет
сократить затраты на действия, не требуемые выбранному классу задач.
Достигаемая экономия зависит от структуры вычислительной нагрузки, условий
исполнения и устройства аппаратной системы. Для её оценки наряду с
характеристиками арифметического узла учитываются затраты памяти и управления.

Проблема памяти также возникла до современных нейросетевых ускорителей.
В статье У.~Вульфа и С.~Макки 1995 года обсуждалось расхождение темпов
развития процессоров и памяти~\cite{wulf1995}. Из-за этого расхождения рост скорости арифметических операций не всегда
приводил к такому же ускорению приложения. В нейросетевых архитектурах
ограничения обмена данными получили новые формы, связанные
с повторным использованием параметров, промежуточных представлений
и частичных результатов вычисления.

\subsection{Глубокое обучение и рост требований к вычислительной аппаратуре}

Работа А.~Крижевского, И.~Суцкевера и Дж.~Хинтона 2012 года показала
результативность большой свёрточной сети в задаче классификации
ImageNet. В реализации использовались графические процессоры
и специально организованные свёрточные вычисления~\cite{krizhevsky2012}.
Эта работа показала связь между увеличением размера обучаемой модели и
доступными аппаратными средствами. Для проведения эксперимента требовалось
организовать выполнение свёрточных операций на выбранных процессорах.

По отношению к ранее разработанным системам анализа изображений
изменились масштаб задачи, объём используемых данных и условия
вычислительного эксперимента. Наряду с качеством модели исследователи учитывали возможность её обучения на
доступной аппаратуре.
Распространение вычислительно сложных моделей усилило значение
организации характерных для них операций и создало предпосылки
для специализированных аппаратных реализаций.

Для обучения и инференса эти условия различаются. Обучение включает
получение градиентов и обновление параметров, тогда как при использовании
модели параметры часто многократно читаются без изменения. Кроме того,
требования к системе инференса включают пропускную способность
и задержку отдельного ответа. При проектировании системы инференса необходимо снижать энергозатраты с учётом
этих требований к скорости и времени ответа.

\subsection{DianNao, Eyeriss и тензорный процессор Google}

В DianNao, представленном Т.~Чэнем и соавторами в 2014 году,
нейросетевые вычисления выполнялись как последовательность
умножения, суммирования и нелинейного преобразования с выделенными
буферами данных~\cite{chen2014}. Рассматривалась небольшая
специализированная схема, а её характеристики оценивались после
проектирования физического размещения в технологическом процессе 65 нм.
Память проектировалась с учётом арифметической части: требовалось подавать
входы и веса, а также сохранять результаты вычислений.

Специализация DianNao определялась структурой нейросетевого
вычисления: составом одновременно используемых данных, организацией
их хранения и последовательностью получения результата нейрона.
Поэтапная редукция произведений задавала порядок суммирования
частичных результатов. Совместное проектирование арифметики
и буферов позволяло учитывать потребности вычисления
при организации передачи операндов.

В Eyeriss Ю.-С.~Чэня, Дж.~Эмера и В.~Сзе 2016 года
исследовалась организация передачи данных~\cite{eyeriss2016}.
Предложенный способ \textit{row-stationary} (RS) был ориентирован
на повторное использование весов, входных значений
и промежуточных результатов в локальной памяти.
В отличие от описания одного арифметического узла,
исследование охватывало передачу разных видов данных
в пределах вычислительного устройства.

Авторы сопоставляли варианты при одинаковых ограничениях
площади и степени параллелизма, используя энергетическую
модель для технологии 65 нм. Для свёрточных слоёв AlexNet
затраты сравниваемых потоков данных превышали затраты RS
в 1,4--2,5 раза в рассмотренных конфигурациях.
Этот расчётный результат показывает, почему организация
обмена стала отдельной задачей проектирования ускорителей.
Публикация связывала описание нового решения
с общей методикой его сравнения с другими вариантами.

Сопоставление DianNao и Eyeriss показывает различие
исследовательских задач середины 2010-х годов.
Наряду с созданием специализированной схемы для нейронных
сетей изучались способы организации обмена внутри такой
схемы. Снижение энергозатрат связывалось с совместным
проектированием памяти и арифметики при сохранении
самой вычислительной задачи.

В публикации Н.~Джуппи и соавторов 2017 года был представлен анализ
первого тензорного процессора Google (Tensor Processing Unit, TPU),
использовавшегося в центрах
обработки данных с 2015 года~\cite{jouppi2017}. Его матричный блок
был предназначен для низкоразрядных целочисленных операций.
Выбор инференса в качестве основной задачи связывал
устройство процессора с требованиями центров обработки
данных и ограничениями времени ответа сервисов.
Публикация включала сравнение со средствами общего назначения,
использовавшимися в той же вычислительной среде.

В TPU систолическая организация матричных операций применялась
в составе промышленной системы инференса. Достигнутые показатели
зависели от подачи весов, ёмкости памяти и свойств приложений.
Показатели изготовленного устройства позволяли оценить результат специализации
для нейросетевых задач в рассмотренных условиях. Загрузка матричного блока
зависела от конфигурации вычислений и требовала отдельного измерения.

Сопоставление трёх проектов показывает, что специализация не сводится
к одному стандартному устройству. Она может выражаться в выделении
нейросетевого конвейера, в выборе потока данных или в крупном регулярном
матричном блоке. Общим становится требование согласовать арифметику,
локальное хранение и передачу операндов. Изменение численной точности также требует согласования этих частей
вычислительного устройства.

\subsection{От производительности к энергетической стоимости задачи}

М.~Горовиц в докладе 2014 года связывал дальнейшее повышение
эффективности с согласованием приложения и аппаратуры~\cite{horowitz2014}.
Стоимость действий зависит от их вида: выполнение арифметической
операции и доставка операндов из разных уровней памяти требуют
разных затрат. Поэтому одинаковое число умножений не означает одинаковую
энергию двух реализаций. Сокращение перемещений может оказаться
для энергозатрат столь же важным, как изменение самого умножителя.

Модель Roofline, предложенная С.~Уильямсом, Э.~Уотерманом
и Д.~Паттерсоном в 2009 году~\cite{roofline2009}, связывает достижимую скорость
арифметических операций с вычислительной мощностью устройства,
пропускной способностью памяти и числом операций на единицу объёма
передаваемых данных. При этом необходимо указать, между какими уровнями памяти передаются данные:
обмен с внешней памятью и обмен внутри кристалла — разные величины.

Эта модель описывает ограничения производительности через соотношение числа
операций и объёма обмена данными независимо от назначения приложения. Один и тот же тип сети
может исполняться в разных режимах при изменении размера пакета,
порядка обработки или аппаратуры. Но Roofline задаёт ограничение
производительности, а не готовое распределение энергетических затрат.
Если скорость ограничена арифметическими узлами, это ещё не означает, что они
потребляют большую часть энергии системы.

Сокращение разрядности меняет арифметические затраты и представление
данных в памяти. Выгода зависит от устранённой работы и дополнительных
действий для преобразования чисел и управления. Так численная точность
становится частью архитектурной задачи.

\subsection{Моделирование ускорителя как исследовательский инструмент}

Развитие специализированной аппаратуры сопровождалось
созданием средств изучения проектов до изготовления кристалла.
Различные варианты памяти и организации исполнения требовали
сравнения, тогда как изготовление каждого варианта было
связано с затратами времени и средств. Архитектурные модели
позволяли исследовать эти решения на этапе проектирования.

В 2019 году А.~Парашар и соавторы представили Timeloop —
средство исследования нейросетевых архитектур и способов
размещения вычислений на них~\cite{timeloop2019}.
В центре работы находилась связь между устройством ускорителя
и организацией выполнения нейросетевых операций.
Оценивался не только отдельный проект, но и набор допустимых
вариантов его использования. Это расширяло возможности
сопоставления решений при общих исходных ограничениях.

В том же году Я.~У, Дж.~Эмер и В.~Сзе предложили Accelergy —
методику архитектурной оценки энергии ускорителей~\cite{accelergy2019}.
Затраты устройства описывались через его компоненты
и выполняемые ими действия. Такое представление позволяло
связывать итоговую оценку с работой памяти, арифметических
и управляющих элементов. Разработка средств оценки
сопровождала разработку самих ускорителей.

Timeloop и Accelergy решали разные, но связанные задачи:
исследование организации вычислений и оценку энергетических
последствий проектных решений. Появление этих средств
характеризует развитие методов исследования аппаратуры.
Описание нового ускорителя дополнялось средствами,
с помощью которых можно было сравнивать его с другими
вариантами и объяснять полученные характеристики.

При использовании архитектурных моделей сохранялось различие
между расчётом проекта и измерением изготовленного устройства.
Также отдельно рассматривалась проверка численных ответов
нейронной сети. Производительность, энергия и качество
решения задачи описывали разные стороны результата.
Их сопоставление связывало разработку алгоритма,
исследование аппаратной организации и экспериментальную проверку.

Для истории энергоэффективных вычислений такие средства
показывают изменение самой исследовательской практики.
Наряду с публикацией характеристик отдельных устройств
появлялись общие способы описания проектов и организации
сравнения. Это позволяло обсуждать не только преимущество
одного ускорителя, но и условия, при которых определённый
способ исполнения давал энергетическую экономию.
